

\section{Le refus du "Direct" : Quand le grec refuse l'oralité}
space{0.3cm}
oindent

L'année 1902 marque une césure fondamentale dans l'histoire de l'enseignement secondaire français,
une rupture qui dépasse largement le cadre administratif pour toucher aux structures mêmes de la
reproduction sociale et culturelle. La réforme portée par le ministre Georges Leygues, aboutissement
des travaux de la commission parlementaire présidée par Alexandre Ribot en 1899, ne se contente pas
de réorganiser les filières du lycée ; elle cristallise les tensions idéologiques d'une France
encore traumatisée par la défaite de 1870 et obsédée par la modernisation de son appareil éducatif
face à la puissance allemande et anglo-saxonne.\cite{I-2-2-ref1} Au cœur de ce dispositif législatif
et réglementaire se joue une bataille silencieuse mais féroce : celle des méthodes d'enseignement
des langues.

Cette période voit l'officialisation spectaculaire de la « méthode directe » pour les langues
vivantes (anglais, allemand), imposant l'oralité, l'immersion et le bannissement de la traduction
comme nouveaux dogmes pédagogiques.\cite{I-2-2-ref2} Pourtant, dans un contraste saisissant,
l'enseignement des langues anciennes, et singulièrement celui du grec ancien, se cabre dans un refus
obstiné de ces nouvelles pratiques. Alors que les classes d'anglais résonnent de conversations et de
descriptions d'images, les classes de grec demeurent le sanctuaire du texte écrit, de l'analyse
grammaticale et de la traduction silencieuse.

Ce rapport se propose d'explorer les racines profondes de ce refus. Il ne s'agit pas ici d'une
simple inertie corporatiste ou d'une inaptitude technique des professeurs de lettres classiques.
L'analyse démontre que le rejet de l'oralité pour le grec ancien obéit à une logique de «
distinction » culturelle et sociale, au sens où l'entendent Pierre Bourdieu et Jean-Claude
Passeron.\cite{I-2-2-ref4} En sanctuarisant le grec comme langue morte, purement graphique et
grammaticale, l'institution scolaire préserve un instrument de sélection des élites, érigeant une
barrière symbolique infranchissable contre l'utilitarisme montant des classes moyennes et la «
commercialisation » du savoir que représente, à leurs yeux, la méthode directe.

Nous analyserons d'abord la rupture épistémologique introduite par la méthode directe dans le
paysage scolaire de 1902, avant d'examiner les raisons spécifiques — linguistiques, historiques et
idéologiques — qui ont conduit à l'exclusion du grec de ce mouvement. Nous démontrerons ensuite
comment ce choix méthodologique a servi de levier sociologique pour maintenir la prééminence de la
filière classique (Section A) dans la hiérarchie des excellences, transformant la maîtrise du grec
écrit en un marqueur de classe indélébile. Enfin, nous aborderons les conséquences de ce clivage,
menant aux contre-réformes des années 1920 et figeant pour longtemps le statut des humanités en
France.



\subsection{Le contexte de la réforme : « coup d'état pédagogique » ou nécessité vitale?}

Pour comprendre la violence des débats autour de la méthode directe, il est impératif de resituer la
réforme de 1902 dans son contexte politique. La IIIe République, solidement installée mais inquiète
de la stagnation démographique et économique de la France, cherche à adapter son élite aux défis du
monde industriel. L'enquête parlementaire de 1899, dirigée par Alexandre Ribot, dresse un constat
alarmant : l'enseignement secondaire, dominé par le monopole du latin et la rhétorique, ne prépare
pas suffisamment aux réalités de la vie moderne.\cite{I-2-2-ref1}

La réforme de 1902 est donc conçue comme une réponse pragmatique. Elle abolit la hiérarchie formelle
entre l'enseignement « classique » et l'enseignement « moderne » en créant un baccalauréat unique,
sanctionnant quatre filières distinctes dans le second cycle :

\begin{itemize}\item \textbf{Section A :} Latin-Grec (l'héritière des humanités traditionnelles).\item \textbf{Section B :} Latin-Langues Vivantes.\item \textbf{Section C :} Latin-Sciences.\item \textbf{Section D :} Sciences-Langues Vivantes (sans latin).\end{itemize}Cette architecture, en apparence égalitaire, masque mal une lutte de prestige. Si les sciences et
les langues vivantes gagnent en volume horaire et en légitimité institutionnelle, elles doivent
prouver leur valeur formatrice. C'est ici qu'intervient la méthodologie. Pour que les langues
vivantes ne soient pas de simples codes utilitaires, mais deviennent de véritables instruments de
culture et d'efficacité, l'administration décide de leur appliquer un traitement de choc : la
méthode directe.



\subsection{La méthode directe : une rupture épistémologique}

L'imposition de la méthode directe par les instructions de 1901 et 1902 (notamment la circulaire du
15 novembre 1901) est vécue comme une révolution. Christian Puren, dans son \textit{Histoire des
méthodologies}, qualifie ce moment de charnière où l'État intervient directement dans la pédagogie
pour imposer une doctrine unique.\cite{I-2-2-ref2}

Qu'est-ce que le « Direct » en 1902? Il se définit par trois refus et trois impératifs, en rupture
totale avec la tradition scolaire :

1. \textbf{Le refus de la traduction :} L'accès au sens ne doit pas passer par la langue maternelle.
On ne traduit pas « \textit{book} » par « livre », on montre l'objet.

2. \textbf{Le refus de la règle \textit{a priori} :} La grammaire doit être inductive. L'élève
découvre la règle par la pratique, il ne l'apprend pas par cœur avant de l'appliquer.

3. \textbf{Le refus du livre (dans un premier temps) :} L'écrit est subordonné à l'oral. L'oreille
et la bouche doivent précéder l'œil.

Les partisans de la méthode directe, inspirés par les réformateurs allemands comme Wilhelm Viëtor
(auteur du pamphlet \textit{Der Sprachunterricht muss umkehren\!} \- \og L'enseignement des langues
doit faire volte-face\!\fg{}) et les pédagogues français comme Charles Schweitzer, défendent l'idée
d'un « bain linguistique ».\cite{I-2-2-ref3} L'objectif est la « possession effective » de la
langue, c'est-à-dire la capacité à communiquer, à comprendre et à agir dans la langue
étrangère.\cite{I-2-2-ref1}

\begin{table}[h!]\small\centering\begin{tabular}{| p{2.3cm} | p{3.5cm} | p{3.5cm} |}\hline\textbf{Aspect Pédagogique} & \textbf{Méthode Traditionnelle (Grammaire-Traduction)} & \textbf{Méthode Directe (1902)} \\ \hline\textbf{Objectif Premier} & Formation logique, accès aux textes littéraires & Communication orale, usage pratique, \og vie\fg{} \\ \hline\textbf{Vecteur d'Apprentissage} & L'écrit (le livre, le dictionnaire) & L'oral (la parole du maître, l'image) \\ \hline\textbf{Rôle de la L1 (Français)} & Central (langue d'explication et de traduction) & Banni (ou strictement limité) \\ \hline\textbf{Statut de la Grammaire} & Explicite, déductive, normative & Implicite, inductive, intuitive \\ \hline\textbf{Support Culturel} & Le passé, les auteurs classiques & Le présent, le \og Realien\fg{} (la vie quotidienne) \\ \hline\end{tabular}\end{table}Cette rupture méthodologique est acceptée, bon gré mal gré, pour l'anglais et l'allemand, car ces
langues sont perçues comme des outils de puissance économique et militaire. Mais dès qu'il s'agit du
grec, la logique s'inverse radicalement.



\subsection{La résistance du grec : le silence de la section a}

Alors que les inspecteurs généraux sillonnent les classes d'anglais pour vérifier que l'on ne parle
plus français, les classes de grec de la Section A restent des îlots de tradition. Aucune circulaire
n'impose l'oralisation du grec ancien. Les instructions de 1902, tout en modernisant les programmes,
confirment implicitement que le grec s'enseigne par la version, le thème et l'analyse
grammaticale.\cite{I-2-2-ref1}

Ce refus n'est pas un oubli. Il est consubstantiel à la définition même de la filière classique.
Dans les débats de la \textit{Revue Universitaire}, les professeurs de langues anciennes observent
avec un mélange de mépris et d'inquiétude l'agitation des classes de langues
vivantes.\cite{I-2-2-ref7} Pour eux, le « Direct » est synonyme de superficialité, de « psittacisme
» (répétition mécanique comme un perroquet) et de renoncement à la haute culture. Le grec refuse
l'oralité parce qu'il refuse de devenir une langue « utile ».

L'analyse des discours de l'époque, notamment ceux de Michel Bréal (professeur au Collège de France)
et des contributeurs de la \textit{Revue Universitaire}, permet de dégager les arguments structurels
qui ont verrouillé l'enseignement du grec dans l'écrit.



\subsection{L'argument ontologique : la nature de la « langue morte »}

L'argument le plus immédiat est celui de la mort de la langue. Michel Bréal, dans ses conférences
\textit{De l'enseignement des langues anciennes} (1891), pose une distinction fondamentale : on
apprend les langues vivantes pour parler aux hommes, on apprend les langues mortes pour parler aux
livres.\cite{I-2-2-ref9}

Pour les pédagogues de 1902, \og parler\fg{} le grec ancien serait une comédie historique. Avec qui
l'élève dialoguerait-il? La reconstitution d'une conversation artificielle sur des sujets de la vie
quotidienne antique (la palestre, l'agora) est jugée indigne de la gravité de la discipline. De
plus, le corpus grec est clos. Contrairement à l'anglais qui évolue et s'invente chaque jour, le
grec est fixé dans une perfection atemporelle, celle du siècle de Périclès. Introduire l'oralité,
c'est introduire le risque de l'improvisation, de la faute, du néologisme barbare, et donc de la
corruption d'un corpus sacré.

Cette vision est renforcée par le fait que la méthode directe repose sur l'intuition et
l'association immédiate mot-objet. Or, l'enseignement du grec vise précisément l'inverse : la
médiation, la réflexion, la distance. Comme le souligne Puren, la méthodologie traditionnelle repose
sur le « paradigme indirect » : comprendre, c'est traduire ; parler (si tant est qu'on le fasse),
c'est réciter ou commenter, jamais échanger.\cite{I-2-2-ref11}



\subsection{Le problème de la prononciation et la « question grecque »}

Un obstacle technique majeur s'oppose à l'oralisation du grec : quelle prononciation adopter? Le
débat est vif et politiquement chargé.

\begin{itemize}\item \textbf{La prononciation érasmienne (restituée) :} Elle est scientifiquement fondée mais artificielle. Elle coupe le grec ancien de ses descendants modernes.\item \textbf{La prononciation moderne (iotacisme) :} Elle est vivante, parlée en Grèce, mais elle introduit des homophonies qui compliquent l'orthographe et, surtout, elle lie le grec ancien au destin de la Grèce moderne.\cite{I-2-2-ref13}\end{itemize}Or, la Grèce de 1902 n'est pas perçue comme un modèle culturel dominant par l'élite française. Elle
est vue à travers le prisme de l'orientalisme ou des questions balkaniques. De plus, en Grèce même,
la guerre linguistique fait rage entre la \textit{Katharevousa} (langue puriste, archaïsante) et la
\textit{Dhimotiki} (langue populaire).\cite{I-2-2-ref13} Importer l'oralité en classe obligerait à
prendre parti dans ces querelles, à reconnaître que le grec est une langue vivante, changeante,
voire \og vulgaire\fg{}.

En refusant l'oralité, l'institution scolaire française neutralise ces tensions. Elle maintient le
grec dans une sphère éthérée, purement textuelle, où la prononciation \og à la française\fg{}
(souvent très éloignée de la réalité antique) suffit puisqu'elle ne sert que de support à la lecture
mentale et à l'analyse grammaticale. Le refus du direct est ici un refus du réel géographique et
historique contemporain.\cite{I-2-2-ref15}



\subsection{La « gymnastique de l'esprit » contre le réflexe conditionné}

L'argumentaire pédagogique des défenseurs du grec repose sur la théorie des facultés. L'esprit se
forme par l'effort, par la résistance que la matière lui oppose. La méthode directe, en cherchant à
créer des automatismes linguistiques (le \textit{pattern drill} avant la lettre), est accusée de
contourner l'intelligence.\cite{I-2-2-ref12}

Pour un professeur de 1902, comprendre une phrase de Thucydide demande une analyse logique
rigoureuse : identification des cas, des modes, des relations syntaxiques complexes. C'est cette «
gymnastique » qui forme le jugement et la raison.\cite{I-2-2-ref16} Si l'élève apprenait le grec \og
naturellement\fg{} comme sa langue maternelle, sans passer par la grammaire explicite, il perdrait
le bénéfice intellectuel de l'exercice.

Le grec refuse l'oralité parce que l'oralité est trop rapide. L'écrit impose la lenteur, le retour
en arrière, la vérification. C'est dans cette lenteur que réside la vertu formatrice de la
discipline, opposée à la fluidité superficielle des échanges verbaux prônés par la méthode directe
en anglais.

Si les arguments pédagogiques et linguistiques forment la superstructure du refus, l'infrastructure
est profondément sociologique. L'analyse de Pierre Bourdieu dans \textit{La Reproduction} (1970)
offre une grille de lecture indispensable pour comprendre les enjeux de 1902 : le grec ancien
fonctionne comme l'outil suprême de la distinction sociale.\cite{I-2-2-ref4}



\subsection{Le « gaspillage ostentatoire » et la gratuité du savoir}

Bourdieu parle de « gaspillage ostentatoire d'apprentissage » pour qualifier l'étude des langues
anciennes.\cite{I-2-2-ref4} Dans une société industrielle qui valorise le temps et l'efficacité,
consacrer des années à étudier une langue qu'on ne parlera jamais est le signe d'une appartenance à
une classe qui n'est pas soumise à l'urgence de la nécessité économique.

La méthode directe en langues vivantes est entachée d'utilitarisme. Elle prépare au commerce, à
l'industrie, au voyage. Elle est \og bourgeoise\fg{} au sens économique, mais pas \og
aristocratique\fg{} au sens culturel. Le grec, lui, est pur parce qu'il est inutile. Refuser de le
rendre \og direct\fg{}, c'est-à-dire efficace et communicatif, c'est préserver sa gratuité. C'est
affirmer que la culture véritable se situe au-delà de l'usage pratique.

La Section A (Latin-Grec) devient ainsi le refuge des héritiers. Pour y réussir, il ne suffit pas
d'avoir de l'intelligence ou de la mémoire (qualités suffisantes pour la méthode directe) ; il faut
posséder un habitus culturel, une familiarité avec l'abstraction littéraire et une maîtrise du
français classique que seule la famille peut transmettre en amont de l'école.\cite{I-2-2-ref4}



\subsection{La barrière du thème et de la version}

La méthode traditionnelle, fondée sur la traduction, agit comme un filtre social puissant.
L'exercice de la version grecque ne teste pas seulement la connaissance du grec, mais surtout la
maîtrise des nuances du français. Il faut savoir choisir entre « craindre », « redouter », «
appréhender » pour traduire un mot grec, en fonction du contexte stylistique.

Cet exercice valorise une culture littéraire patrimoniale. À l'inverse, la méthode directe, qui
prône l'expression spontanée dans la langue cible, est potentiellement plus démocratique (ou du
moins méritocratique) : un élève doué peut acquérir une bonne compétence orale en classe, même s'il
ne possède pas la bibliothèque de la Pléiade chez lui. En maintenant le grec dans la méthode
grammaire-traduction, l'institution scolaire s'assure que la sélection se fait sur des critères
culturels implicites, favorisant les enfants de la bourgeoisie lettrée.\cite{I-2-2-ref19}



\subsection{La stratégie de défense des humanités contre la « masse »}

La réforme de 1902 a ouvert les vannes de l'enseignement secondaire. Les effectifs augmentent, de
nouvelles couches sociales accèdent au lycée (notamment via les sections modernes). Face à cette \og
massification\fg{} relative, les humanités classiques doivent se distinguer pour
survivre.\cite{I-2-2-ref21}

Le refus de l'oralité est une stratégie de clôture. En rendant l'apprentissage du grec aride,
difficile, austère, on en fait une épreuve initiatique. Ceux qui survivoient à la grammaire grecque
sans méthode directe prouvent leur valeur morale (la discipline, l'effort) et leur appartenance à
l'élite. Si le grec devenait \og amusant\fg{} ou \og facile\fg{} grâce à des méthodes actives, il
perdrait sa valeur de tri. Comme le note Bourdieu, « ce que l'on cherche à faire croire aux
dépossédés, c'est qu'ils n'ont pas les dons nécessaires », alors qu'ils n'ont simplement pas les
clés du code culturel.\cite{I-2-2-ref4} Le grec enseigné traditionnellement \textit{est} ce code.

Le refus du direct en grec ne s'est pas fait sans heurts. Il s'inscrit dans une guerre de tranchées
qui a duré deux décennies, culminant avec la contre-réforme de 1923.



\subsection{L'anglais et l'allemand comme contre-modèles repoussoirs}

L'application de la méthode directe en anglais et en allemand a servi de repoussoir absolu pour les
défenseurs du grec. Les revues pédagogiques de l'époque, comme \textit{Les Langues Modernes} (fondée
par l'APLV) ou la \textit{Revue de l'Enseignement des Langues Vivantes} (RELV), documentent cette
tension.\cite{I-2-2-ref6}

Les partisans des langues anciennes, souvent soutenus par des figures de la \textit{Revue
Universitaire}, caricaturent la méthode directe. Ils décrivent des classes transformées en foires,
où l'on gesticule pour se faire comprendre, où l'on parle un « petit nègre » (terme raciste de
l'époque désignant un langage simplifié) et où la haute littérature est abandonnée au profit de la
lecture des journaux ou des catalogues commerciaux.\cite{I-2-2-ref7}

Cette caricature permet de justifier par contraste la noblesse du grec. Si l'allemand est la langue
des \og choses\fg{} (Realien), le grec reste la langue des \og idées\fg{}. Le refus de l'oralité est
une manière de ne pas contaminer l'esprit de l'élève par le matérialisme ambiant.



\subsection{La « crise du français » : l'argument ultime}

Dès 1910, une panique morale s'empare des milieux éducatifs : la « Crise du français ». On déplore
que les bacheliers ne sachent plus écrire, que l'orthographe se perde, que la syntaxe se
relâche.\cite{I-2-2-ref24}

Les coupables sont vite désignés : la méthode directe et la réduction de la part du latin/grec. Les
conservateurs, regroupés dans des ligues comme la \textit{Ligue pour la culture française} (soutenue
par des académiciens comme Jean Richepin ou des philosophes comme Henri Bergson, bien que ce dernier
soit plus nuancé), développent l'argument suivant :

1. La méthode directe a supprimé la traduction (thème/version) en langues vivantes.

2. Les élèves ne pratiquent plus assez la comparaison entre les langues et le français.

3. Seul le grec (et le latin), par sa méthode traditionnelle de traduction mot à mot et d'analyse,
continue de former à la rigueur de la langue française.

Ainsi, le refus de l'oralité en grec devient un acte de patriotisme linguistique. On ne parle pas
grec pour mieux écrire français. Cet argument sera le fer de lance de la réaction politique des
années 1920.\cite{I-2-2-ref26}



\subsection{La réforme bérard de 1923 : le triomphe de la distinction}

L'aboutissement politique de ce refus est la réforme portée par Léon Bérard en 1923. Ministre de
l'Instruction publique, Bérard supprime le cycle moderne sans latin en 6e et 5e, rétablissant
l'obligation des humanités classiques pour tous les élèves du lycée (avant que le Cartel des gauches
ne revienne dessus en 1924).\cite{I-2-2-ref7}

Les arguments de Bérard sont explicitement liés à la distinction culturelle. Il affirme que
l'enseignement \og moderne\fg{} (scientifique et langues vivantes par méthode directe) a échoué à
former une véritable élite. Seules les humanités classiques, avec leur méthode lente et difficile,
garantissent la \og culture désintéressée\fg{} nécessaire aux dirigeants. En 1923, l'État français
valide officiellement la thèse selon laquelle la méthode traditionnelle (refus du direct, primauté
de l'écrit et du passé) est supérieure à la méthode moderne pour la formation de
l'esprit.\cite{I-2-2-ref29}



\subsection{L'exception française? comparaison avec l'allemagne et l'angleterre}

Il est éclairant de comparer la situation française avec ses voisins.

\begin{itemize}\item \textbf{Allemagne :} Patrie de la philologie et de la méthode directe (Viëtor), l'Allemagne connaît aussi une tension entre \textit{Gymnasium} (classique) et \textit{Realschule}. Cependant, la tradition philologique allemande (Altertumswissenschaft) est si puissante qu'elle intègre plus facilement une approche scientifique et vivante des textes, sans nécessairement passer par le rejet total de l'oralité pédagogique, bien que la traduction reste centrale.\cite{I-2-2-ref6}\item \textbf{Angleterre :} La méthode directe y connaît un grand succès pour le français, mais les \textit{Public Schools} (Eton, Harrow) maintiennent un enseignement du grec et du latin extrêmement traditionnel (versification, traduction), fonctionnant là aussi comme un marqueur de classe sociale très fort (le \og Classicist\fg{} est au sommet de la hiérarchie sociale britannique).\end{itemize}La spécificité française réside dans la politisation extrême du débat (l'école est le champ de
bataille de la République) et dans la rigidité de la coupure entre méthode directe (réservée au
vivant/utile) et méthode traditionnelle (réservée au mort/culturel).



\subsection{L'héritage du refus : un siècle de « grec mort »}

Le refus de 1902 a eu des conséquences durables. Pendant près d'un siècle, l'enseignement du grec en
France est resté figé dans le modèle de la grammaire-traduction. Les tentatives d'introduire de
l'oralité ou de la méthode active ont été marginalisées jusqu'aux années 1980-1990 (travaux de Jean-
Pierre Vernant ou renouveau didactique).\cite{I-2-2-ref31}

Ce n'est que très récemment, avec le succès paradoxal de méthodes comme \textit{Polis} (Christophe
Rico) ou la méthode Ørberg (naturelle), que l'idée de parler le grec ancien (koinè) refait surface,
souvent en dehors de l'institution scolaire classique ou dans des cercles
passionnés.\cite{I-2-2-ref32} Mais en 1902, cette voie était barrée par la nécessité sociologique de
la distinction.

L'analyse de la conjoncture de 1902 révèle que le refus de la méthode directe pour le grec ancien
n'était pas un accident de l'histoire, ni une simple résistance au changement. C'était une opération
de sauvegarde symbolique.

Face à la montée des valeurs bourgeoises utilitaristes, incarnées par la méthode directe en anglais
et en allemand, l'université française et les élites républicaines ont érigé le grec ancien en
citadelle de la « culture pure ».

\begin{itemize}\item \textbf{Au nom de la pédagogie}, ils ont opposé la gymnastique consciente de l'esprit à l'automatisme inconscient du réflexe.\item \textbf{Au nom de la linguistique}, ils ont opposé la perfection close de la langue morte à la corruption vivante de l'usage.\item \textbf{Au nom de la sociologie} (inconsciemment ou non), ils ont opposé la distinction de l'héritier capable de perdre son temps à traduire Platon, à la vulgarité du parvenu pressé de parler anglais pour vendre ses marchandises.\end{itemize}Le grec a refusé l'oralité pour rester une langue de classe, une langue de caste, une langue de
papier dont la difficulté même était la garantie de sa valeur. En 1902, on a choisi de ne pas parler
le grec pour mieux le faire résonner comme l'instrument silencieux de la domination culturelle.

Ce tableau synthétise les oppositions structurelles qui ont rendu impossible l'adoption du \og
Direct\fg{} en grec.

\begin{table}[h!]\small\centering\begin{tabular}{| p{2cm} | p{2.2cm} | p{2.2cm} | p{2.2cm} |}\hline\textbf{Critère} & \textbf{Monde des Langues Vivantes (Réforme 1902)} & \textbf{Monde du Grec Ancien (Section A)} & \textbf{Enjeu Sociologique (Bourdieu)} \\ \hline\textbf{Méthode} & Directe (Immersive, Orale) & Traditionnelle (Grammaire-Traduction) & Efficacité vs Distinction \\ \hline\textbf{Finalité} & Utilité sociale (Commerce, Guerre, Science) & Culture désintéressée (Humanités) & Capital Économique vs Capital Culturel \\ \hline\textbf{Exercice Type} & Conversation, Description d'image & Version, Thème, Analyse logique & Compétence technique vs Habitus bourgeois \\ \hline\textbf{Rapport au Temps} & Rapidité, Réflexe, Immédiateté & Lenteur, Réflexion, Médiation & Le temps comme ressource vs Le temps comme luxe \\ \hline\textbf{Public Cible} & Classes moyennes montantes (Sections B, D) & Élite traditionnelle et héritiers (Section A) & Ouverture démocratique vs Reproduction des élites \\ \hline\textbf{Statut de l'Erreur} & Tolérée si la communication passe & Sanctionnée comme faute logique ou de goût & Pragmatique vs Normative \\ \hline\end{tabular}\end{table}